% Discussion -- cause of the error rise. Numbers: results/cause_checks.csv (pipeline/cause_checks.py)
\paragraph{What the error rise is not caused by.}
We do not identify the mechanism behind the larger errors at high multireference character, but three checks at the Transition1x transition states narrow it down (SI, Section~S6). First, the rise is not caused by the broken-symmetry labels in OMol25: MACE-T1x, which was trained only on restricted labels, shows the same factor of about two between low and high $N_\mathrm{FOD}$ as the OMol25 models (Section~\ref{sec:labelcheck}), and at equal $N_\mathrm{FOD}$ its errors hardly differ between the spin classes (factor 1.06). Second, it is not caused by inconsistent labels. Along the Transition1x paths of 275 reactions, the OMol25 labels switch between solutions more often at high $N_\mathrm{FOD}$, but these switches do not predict the model error once $N_\mathrm{FOD}$ is accounted for (partial rank correlation 0.05--0.10). Third, the error rises equally with a smaller HOMO--LUMO gap, which is closely correlated with $N_\mathrm{FOD}$ (rank correlation $-0.81$), so both point to near-degenerate orbitals. The larger errors therefore appear to come from the electronic structure itself, independent of how the reference treats it. How near-degeneracy makes the surface harder to learn, for example through a less local dependence of the energy on the geometry or through too few such structures in the training data, remains open.
